\begin{Exercise}[label={ex:capocannoniere}]
		Un campionato di calcio \`e composto da venti squadre. Ciascuna squadra ha un nome (lungo al pi\'u 30 caratteri), un numero di punti accumulati durante la stagione e al più 23 calciatori. Ciascun calciatore ha un cognome (lungo al pi\'u 20 caratteri), un'et\'a, un numero di maglia, un numero di presenze in stagione ed un numero di gol realizzati in stagione.
		\begin{enumerate}
			\item Specificare i tipi Campionato, Squadra e Calciatore.
			\item Viene data una variabile di nome serie\_a di tipo Campionato, gi\'a inizializzata e riempita dall'utente con tutte le squadre del campionato. Scrivere un frammento di codice che stampi il cognome, il numero di gol realizzati e la squadra di appartenenza del calciatore che ha realizzato pi\'u gol nel campionato.
		\end{enumerate}
\end{Exercise}

\begin{Answer}[ref={ex:capocannoniere}]
	\begin{enumerate}
		\item \lstinputlisting[style=cstylesmall]{src/campionato.c}
		
		\item \lstinputlisting[style=cstylesmall]{src/capocannoniere.c}
	\end{enumerate}
\end{Answer}